\documentclass[12pt]{article}
\usepackage{amsmath}
\usepackage{geometry}
\geometry{margin=1in}

\begin{document}

\title{Elastic Fields of a Screw Dislocation}
\author{}
\date{}
\maketitle

\section*{1. Displacement Field (Screw Dislocation)}

For a screw dislocation with Burgers vector $b$ along the $z$-axis:

\[
u_z = \frac{b}{2\pi}\tan^{-1}\!\left(\frac{y}{x}\right),
\qquad
u_x = 0, \qquad u_y = 0.
\]

Let
\[
\theta = \tan^{-1}\!\left(\frac{y}{x}\right), \qquad
r = \sqrt{x^2+y^2}.
\]


\section*{2. Strain Field Derivations}

All normal strains vanish:

\[
\varepsilon_{xx} = 0, \qquad
\varepsilon_{yy} = 0, \qquad
\varepsilon_{zz} = 0.
\]

Shear strains are nonzero only in components involving $z$:

\[
\varepsilon_{xz}
= \frac{1}{2}\frac{\partial u_z}{\partial x}
= -\frac{b}{4\pi r}\sin\theta,
\]

\[
\varepsilon_{yz}
= \frac{1}{2}\frac{\partial u_z}{\partial y}
= \frac{b}{4\pi r}\cos\theta.
\]

\[
\varepsilon_{xy} = 0.
\]

Trace:

\[
\varepsilon_{kk} = 0.
\]

\section*{3. Hooke's Law (Isotropic Elasticity)}

Since the only nonzero strains are shear strains and the trace is zero:

\[
\sigma_{xz} = 2G\varepsilon_{xz},
\qquad
\sigma_{yz} = 2G\varepsilon_{yz},
\qquad
\sigma_{xx} = \sigma_{yy} = \sigma_{zz} = \sigma_{xy} = 0.
\]

\newpage
\section*{4. Stress Field Derivations}

\subsection*{$\sigma_{xz}$}

\[
\sigma_{xz}
= 2G\varepsilon_{xz}
= 2G\left(-\frac{b}{4\pi r}\sin\theta\right),
\]

\[
\boxed{
\sigma_{xz}
= -\frac{Gb}{2\pi r}\sin\theta
}
\]

\subsection*{$\sigma_{yz}$}

\[
\sigma_{yz}
= 2G\varepsilon_{yz}
= 2G\left(\frac{b}{4\pi r}\cos\theta\right),
\]

\[
\boxed{
\sigma_{yz}
= \frac{Gb}{2\pi r}\cos\theta
}
\]


\section*{5. Final Classical Stress Fields (Screw Dislocation)}

\[
\boxed{
\sigma_{xz}
= -\frac{Gb}{2\pi r}\sin\theta
}
\]

\[
\boxed{
\sigma_{yz}
= \frac{Gb}{2\pi r}\cos\theta
}
\]

\[
\sigma_{xx}=\sigma_{yy}=\sigma_{zz}=\sigma_{xy}=0.
\]

These are the standard textbook stresses of a screw dislocation.

\end{document}
